
\SGASubsection{4.2}{Radicial Groups}
\SGARefTarget{sga3:xvii:subsection:4:2}

\SGAStatement{\SGAInlineRefTarget{sga3:xvii:proposition:4-2-1}Proposition 4.2.1}{Let \(G\) be a radicial algebraic
group defined over a field \(k\) of characteristic \(p>0\).  The
following conditions are equivalent:
\begin{enumerate}
\item[(i)] \SGARefTarget{sga3:xvii:proposition:4:2:1:item:i}
  \(G\) is unipotent.
\item[(ii)] \SGARefTarget{sga3:xvii:proposition:4:2:1:item:ii}
  \(G\) has a central composition series whose successive
  quotients are isomorphic to \(\boldsymbol\alpha_p\).
\item[(iii)] \SGARefTarget{sga3:xvii:proposition:4:2:1:item:iii}
  \(G\) has a central, characteristic composition series whose
  successive quotients are isomorphic to
  \((\boldsymbol\alpha_p)^r\).
\item[(iv)] \SGARefTarget{sga3:xvii:proposition:4:2:1:item:iv}
  \(G_{\overline{k}}\) contains no subgroup isomorphic to
  \(\boldsymbol\mu_p\).
\item[(v)] \SGARefTarget{sga3:xvii:proposition:4:2:1:item:v}
  \(\mathfrak g=\operatorname{Lie}G\) is a unipotent
  \(p\)-Lie algebra (\SGARefLink{sga3:xvii:definition:3:6}{3.6}).
\end{enumerate}}
\SGARefTarget{sga3:xvii:proposition:4:2:1}

\SGAProof{Proof}
The implications
\SGARefLink{sga3:xvii:proposition:4:2:1:item:iii}{(iii)}\(\Rightarrow\)\SGARefLink{sga3:xvii:proposition:4:2:1:item:ii}{(ii)}\(\Rightarrow\)\SGARefLink{sga3:xvii:proposition:4:2:1:item:i}{(i)} are clear;
\SGARefLink{sga3:xvii:proposition:4:2:1:item:i}{(i)}\(\Rightarrow\)\SGARefLink{sga3:xvii:proposition:4:2:1:item:iii}{(iii)} is \SGARefLink{sga3:xvii:lemma:3:9:item:i}{3.9(i)}, and
\SGARefLink{sga3:xvii:proposition:4:2:1:item:i}{(i)}\(\Rightarrow\)\SGARefLink{sga3:xvii:proposition:4:2:1:item:iv}{(iv)} follows from \SGARefLink{sga3:xvii:proposition:2:4:item:i}{2.4(i)}.

We shall need the following lemma on abelian \(p\)-Lie algebras.

\SGAStatement{\SGAInlineRefTarget{sga3:xvii:lemma:4-2-2}Lemma 4.2.2}{Let \(\mathfrak g\) be an abelian,
finite-dimensional \(p\)-Lie algebra over a perfect field \(k\).
Then \(\mathfrak g\) can be written uniquely as a direct sum
\[
  \mathfrak g=\mathfrak r\oplus\mathfrak u
\]
of a \(p\)-Lie subalgebra \(\mathfrak r\) on which the \(p\)-th
power is bijective and a unipotent \(p\)-Lie subalgebra
\(\mathfrak u\) (\SGARefLink{sga3:xvii:definition:3:6}{3.6}).  The algebra \(\mathfrak r\) is called the
\emph{reductive part} of \(\mathfrak g\), and \(\mathfrak u\) its
\emph{unipotent part}.  The formation of \(\mathfrak r\) and
\(\mathfrak u\) is compatible with extension of the field \(k\).
If, moreover, \(k\) is algebraically closed, then
\(\mathfrak r\) has a basis \(e_i\) such that
\[
  e_i^{(p)}=e_i.
\]}
\SGARefTarget{sga3:xvii:lemma:4:2:2}

The proof of this lemma is easy and is left to the reader; see
Bourbaki, \emph{Groupes et algèbres de Lie}, Chapter I, \S1,
Exercise 23.  We merely note that \(\mathfrak u\) is the kernel of
a suitable iterate of the map \(X\mapsto X^{(p)}\), while
\(\mathfrak r\) is the image of that same iterate.

We now prove \SGARefLink{sga3:xvii:proposition:4:2:1:item:iv}{(iv)}\(\Rightarrow\)\SGARefLink{sga3:xvii:proposition:4:2:1:item:v}{(v)}.  After extending the base
field if necessary, we may suppose that \(k\) is algebraically
closed.  Let \(X\in\mathfrak g\), and let \(\mathfrak h\) be the
\(p\)-Lie subalgebra of \(\mathfrak g\) generated by \(X\).  The
algebra \(\mathfrak h\) is plainly commutative, and its reductive
part (\SGARefLink{sga3:xvii:lemma:4-2-2}{4.2.2}) is zero.  Otherwise, by the cited result,
\(\mathfrak h\) would contain an element \(Y\ne0\) such that
\(Y^{(p)}=Y\), and consequently, by Appendix II, \SGARefLink{sga3:xvii:proposition:2:1}{2.1} and \SGARefLink{sga3:xvii:proposition:2:2}{2.2},
\(G\) would contain a subgroup isomorphic to
\(\boldsymbol\mu_p\), contrary to the hypothesis.  Thus
\(\mathfrak h\), and hence \(\mathfrak g\), is a unipotent
\(p\)-Lie algebra.

It remains to prove \SGARefLink{sga3:xvii:proposition:4:2:1:item:v}{(v)}\(\Rightarrow\)\SGARefLink{sga3:xvii:proposition:4:2:1:item:i}{(i)}, the least immediate
implication in \SGARefLink{sga3:xvii:proposition:4-2-1}{4.2.1}.

\medskip\noindent\textit{(a) The case in which \(G\) has height
\(1\)} (Exposé VII\textsubscript{A}, \SGARefLink{sga3:VIIA:4.1.3}{4.1.3}).
Since \(G\) is radicial, it is affine and hence isomorphic to an
algebraic subgroup of a linear group \(\operatorname{GL}(V)\)
(Exposé VI\textsubscript{B}, \S11).  This embedding identifies
\(\mathfrak g\) with a \(p\)-Lie subalgebra of
\(\operatorname{End}(V)\), and the \(p\)-th power of \(X\) in
\(\mathfrak g\) agrees with the \(p\)-th power of the endomorphism
\(X\) (Exposé VII\textsubscript{A}, \SGARefLink{sga3:VIIA:6.4.4}{6.4.4}).  Since
\(\mathfrak g\) is unipotent by hypothesis, it is therefore an
algebra of nilpotent endomorphisms of \(V\).  By Engel's theorem
(Bourbaki, \emph{Groupes et algèbres de Lie}, Chapter I, \S4,
Theorem 1), it is a Lie subalgebra of the Lie algebra
\(\mathfrak h\) of the group of strictly upper-triangular matrices
\(\operatorname{Trigstr}(n)\), relative to a suitable basis of
\(V\).  Since \(G\) has height \(1\), Appendix II, \SGARefLink{sga3:ii:subsection:2-2}{2.2} then shows
that \(G\) itself is an algebraic subgroup of
\(\operatorname{Trigstr}(n)\), and hence is unipotent by
\SGARefLink{sga3:xvii:theorem:3:5:item:v}{3.5(v)}\(\Rightarrow\)\SGARefLink{sga3:xvii:proposition:4:2:1:item:i}{(i)}.

